\section{Every centered moment of two even harmonic tails}
\label{tail:section}

The latest harmonic-parity report asks for exact centered moments of
\((\zeta(2i)-H_n^{(2i)})(\zeta(2j)-H_n^{(2j)})\), beginning with
\((i,j)=(1,2)\). Its all-degree evaluation covers only the trigamma
square. The following results answer the question for every pair of
even orders and every nonpositive even spectral argument. They also
give a convergent generating function in classical polylogarithms.
Products of four or more tails, and products with additional powers
of \(H_n\), are outside the theorem.

\subsection{Definitions, continuation, and an ordinary convergent sum}

Throughout this section, \(p,q\ge2\) are even integers, \(W=p+q\), and
\[
 A_p(n)=\zeta(p)-H_n^{(p)}=\zeta(p,n+1)
       =\frac{\psi^{(p-1)}(n+1)}{(p-1)!},
 \qquad x_n=n+\tfrac12.
\]
The last equality uses that \(p\) is even. Define
\begin{equation}\label{tail:D}
 D_{p,q}(s)=\sum_{n=0}^{\infty}\frac{A_p(n)A_q(n)}{x_n^s},
 \qquad \Re s>3-p-q.
\end{equation}
The defining series converges normally in this half-plane. A value
outside it always refers to meromorphic continuation unless it is
written as an explicitly subtracted convergent sum.

Set
\begin{equation}\label{tail:asymptotic-coefficients}
 c_{p,j}=\frac1{p-1}\binom{p+2j-2}{p-2}B_{2j}(1/2),
 \qquad
 k_{p,q,\ell}=\sum_{j=0}^{\ell}c_{p,j}c_{q,\ell-j}.
\end{equation}
Thus \(c_{p,0}=1/(p-1)\). The standard Hurwitz-zeta expansion gives
\begin{equation}\label{tail:centered-expansion}
 A_p(n)\sim\sum_{j\ge0}c_{p,j}x_n^{1-p-2j},
 \qquad
 A_p(n)A_q(n)\sim\sum_{\ell\ge0}
                  k_{p,q,\ell}x_n^{2-W-2\ell}.
\end{equation}
Every exponent in the product is even. The expansions and their
fixed-order remainder estimates follow either from the classical
Hurwitz formulas or directly from the Laplace integral used below
\cite{DLMF}.

\begin{proposition}[Continuation and ordinary summation]
\label{tail:ordinary}
The function \(D_{p,q}\) continues meromorphically to the plane. Its
only possible poles are simple and lie at
\[
 3-p-q,\ 1-p-q,\ -1-p-q,\ldots.
\]
In particular it is regular at every \(s=-2r\), \(r\ge0\).
Let
\[
 K_r=\max\{0,r-W/2+2\}.
\]
Then
\begin{equation}\label{tail:ordinary-formula}
 D_{p,q}(-2r)=\sum_{n=0}^{\infty}
 \left\{x_n^{2r}A_p(n)A_q(n)
       -\sum_{\ell=0}^{K_r-1}
         k_{p,q,\ell}x_n^{2r+2-W-2\ell}\right\}.
\end{equation}
An empty inner sum is zero. Every bracket is \(O(n^{-2})\) or
smaller, so the right side is an ordinary absolutely convergent sum.
\end{proposition}

\begin{proof}
Subtract any fixed number of terms from the product expansion in
\eqref{tail:centered-expansion}. The remainder series extends to a
larger left half-plane; the removed terms sum to
\(k_{p,q,\ell}\zeta(s+W-2+2\ell,1/2)\). Increasing the number of
subtractions proves the continuation and pole list. At \(s=-2r\),
the nondecaying terms are exactly those with \(0\le\ell<K_r\).
Their continued sums vanish, because
\(\zeta(-2m,1/2)=0\) for every integer \(m\ge0\). The remaining
sum converges absolutely and gives \eqref{tail:ordinary-formula}.
\end{proof}

The half-shift is essential to this argument: it makes both the
asymptotic powers and the polynomial subtraction compatible with the
trivial zeros of the same Hurwitz zeta function.

\subsection{A finite formula at every moment order}

Our double-zeta convention is
\[
 \zeta(a,b)=\sum_{\ell>k\ge1}\ell^{-a}k^{-b},
 \qquad a\ge2,\quad b\ge1.
\]
Only convergent double zeta values occur in the formulas below.
For a positive odd integer \(d\), define the ordered contribution
\begin{equation}\label{tail:ordered-block}
 U_{p,q}(d)=
 \begin{cases}
 \displaystyle
 \zeta(p,q-d)+\frac12\zeta(W-d),&d<q,\\[6pt]
 \displaystyle
 \sum_{h=1}^{\min\{(d-q+1)/2,\,p/2-1\}}
 \frac{\binom{d-q+1}{2h}}{d-q+1}
 B_{d-q+1-2h}\zeta(p-2h),&d>q.
 \end{cases}
\end{equation}
The two cases exhaust the possibilities since \(q\) is even.
The upper limit in the second line is an integer; the sum may be
empty. Define also the rational contribution
\begin{equation}\label{tail:rational-block}
 R_{p,q}(d)=
 \begin{cases}
 0,&d<W-1,\\[3pt]
 \displaystyle
 -\frac{B_{d-W+1}}2
 \left\{\frac1p\binom{d-q}{p-1}
          +\frac1q\binom{d-p}{q-1}\right\},&d\ge W-1,
 \end{cases}
\end{equation}
and
\begin{equation}\label{tail:block}
 C_{p,q}(d)=U_{p,q}(d)+U_{q,p}(d)+R_{p,q}(d).
\end{equation}
These expressions use only finite sums of explicitly specified
Bernoulli numbers and ordinary or double zeta values.

\begin{theorem}[All even orders and all centered moments]
\label{tail:all-moments}
For even \(p,q\ge2\) and every integer \(r\ge0\),
\begin{equation}\label{tail:moment-formula}
 \boxed{
 D_{p,q}(-2r)=
 \sum_{j=0}^{r}\frac{\binom{2r}{2j}}{2j+1}
 B_{2r-2j}(1/2)\,C_{p,q}(2j+1).}
\end{equation}
In particular this formula evaluates the ordinary sums in
\eqref{tail:ordinary-formula}.
\end{theorem}

\begin{proof}
Write \(\operatorname{CT}_N\) for the coefficient of
\(N^0(\log N)^0\) in an asymptotic expansion at positive infinity.
This coefficient will only be taken after exact finite summation.
Let
\begin{equation}\label{tail:S}
 S_r(k)=\sum_{n=0}^{k-1}(n+1/2)^{2r}
       =\frac{B_{2r+1}(k+1/2)}{2r+1}
       =\sum_{j=0}^{r}
         \frac{\binom{2r}{2j}}{2j+1}
         B_{2r-2j}(1/2)k^{2j+1}.
\end{equation}
This is an odd polynomial in \(k\), with zero constant coefficient.
Consequently every polynomial partial sum of a subtracted even power
in \eqref{tail:ordinary-formula} has zero constant coefficient in
\(N\). It follows that
\begin{equation}\label{tail:CT-moment}
 D_{p,q}(-2r)=\operatorname{CT}_N
              \sum_{n=0}^{N-1}x_n^{2r}A_p(n)A_q(n).
\end{equation}

The identity \(A_p(k-1)=A_p(k)+k^{-p}\) yields finite summation by
parts:
\begin{align}\label{tail:finite-summation}
 \sum_{n=0}^{N-1}x_n^{2r}A_p(n)A_q(n)
 ={}&S_r(N)A_p(N)A_q(N)\notag\\
 &+\sum_{k=1}^{N}S_r(k)
  \left\{\frac{A_p(k)}{k^q}+\frac{A_q(k)}{k^p}
                     +\frac1{k^{p+q}}\right\}.
\end{align}
For a positive odd \(d\), the coefficient block to evaluate is
\begin{equation}\label{tail:E}
 E_{p,q,d}(N)=N^dA_p(N)A_q(N)
 +\sum_{k=1}^{N}k^{d-q}A_p(k)
 +\sum_{k=1}^{N}k^{d-p}A_q(k)
 +\sum_{k=1}^{N}k^{d-W}.
\end{equation}
We prove that its constant coefficient is \(C_{p,q}(d)\).

First suppose \(d<q\). The corresponding tail sum converges, and
its limit is exactly
\begin{equation}\label{tail:negative-power}
 \sum_{k\ge1}k^{d-q}A_p(k)=\zeta(p,q-d).
\end{equation}
No continued or divergent double sum is used. If \(m=d-q\ge1\),
then \(m\) is odd. Put
\[
 F_m(N)=\sum_{k=1}^{N}k^m
       =\frac{B_{m+1}(N+1)-B_{m+1}}{m+1}.
\]
Interchanging a finite sum gives the exact identity
\begin{equation}\label{tail:Faulhaber-change}
 \sum_{k=1}^{N}k^mA_p(k)
 =A_p(N)F_m(N)
   +\sum_{\ell=1}^{N}\frac{F_m(\ell-1)}{\ell^p}.
\end{equation}
Since
\[
 F_m(\ell-1)=\frac1{m+1}
    \sum_{a=1}^{m+1}\binom{m+1}{a}B_{m+1-a}\ell^a,
\]
the constant term of the last sum is obtained directly from finite
power sums. Exponents \(a-p\ge0\) give zero constant; exponent
\(-1\) gives \(\gamma\); smaller exponents give \(\zeta(p-a)\).
Because \(m\) is odd, the nonzero Bernoulli coefficients have
\(a\) even, except for the half-leading term \(a=m\), whose
coefficient is \(-1/2\). Thus the even values contribute
\begin{equation}\label{tail:even-zeta-block}
 \sum_{h=1}^{\min\{(m+1)/2,p/2-1\}}
 \frac{\binom{m+1}{2h}}{m+1}B_{m+1-2h}\zeta(p-2h).
\end{equation}
If \(d<W-1\), the half-leading term contributes
\(-\zeta(W-d)/2\). At \(d=W-1\), it instead contributes
\(-\gamma/2\). If \(d>W-1\), it contributes zero.

The final power sum in \eqref{tail:E} contributes
\(\zeta(W-d)\) for \(d<W-1\), \(\gamma\) for \(d=W-1\), and
zero for \(d>W-1\). Combining it with the two half-leading terms
has a useful consequence. Each ordered block with \(d<q\) retains
one half of \(\zeta(W-d)\), while every block with \(d>q\)
retains none. At \(d=W-1\), the two terms \(-\gamma/2\) cancel
the final \(+\gamma\) exactly. Equations
\eqref{tail:negative-power} and \eqref{tail:even-zeta-block}
therefore give precisely \(U_{p,q}(d)+U_{q,p}(d)\).

It remains to compute the rational boundary contributions. The
uncentered expansion is
\begin{equation}\label{tail:uncentered}
 A_p(N)\sim\frac{N^{1-p}}{p-1}-\frac12N^{-p}
 +\sum_{j\ge1}\frac1{p-1}
       \binom{p+2j-2}{p-2}B_{2j}N^{1-p-2j}.
\end{equation}
For any odd integer \(k\), define the inverse-power coefficient by
\[
 u_p(k)=
 \begin{cases}
 0,&k<p-1,\\[2pt]
 \displaystyle\frac1{p-1}\binom{k-1}{p-2}B_{k-p+1},&k\ge p-1.
 \end{cases}
\]
The only nonzero even inverse-power coefficient in
\eqref{tail:uncentered} is \(-1/2\), at exponent \(p\).
In the Faulhaber polynomial \(F_m(N)\), all powers are even
except for its half-leading term \(N^m/2\). Pairing the possible
exponents therefore gives, for odd \(m\ge1\),
\begin{equation}\label{tail:single-boundary}
 \operatorname{CT}_N\{A_p(N)F_m(N)\}
       =\frac{p-m}{2p}\,u_p(m).
\end{equation}
Indeed the two potential contributions are
\(u_p(m)/2\) and
\(-\binom{m+1}{p}B_{m+1-p}/(2(m+1))\); the latter equals
\(-m u_p(m)/(2p)\). If \(m<p-1\), both are zero.

The odd coefficient of the product in the first term of
\eqref{tail:E} is
\[
 \operatorname{CT}_N\{N^dA_p(N)A_q(N)\}
      =-\frac12\{u_p(d-q)+u_q(d-p)\},
\]
where a term below its support is zero. Adding this to the two
contributions \eqref{tail:single-boundary} gives
\[
 -\frac{d-q}{2p}u_p(d-q)
 -\frac{d-p}{2q}u_q(d-p).
\]
It vanishes for \(d<W-1\). Otherwise, the binomial identity
\[
 \frac{d-q}{p-1}\binom{d-q-1}{p-2}
          =\binom{d-q}{p-1}
\]
reduces it to \(R_{p,q}(d)\) in \eqref{tail:rational-block}.
We have proved \(\operatorname{CT}_NE_{p,q,d}(N)=C_{p,q}(d)\).
Combining \eqref{tail:S}, \eqref{tail:CT-moment}, and
\eqref{tail:finite-summation} proves the theorem.
\end{proof}

\subsection{Reduction to a fixed set of ordinary zeta coordinates}

There is no unresolved depth-two constant in the preceding theorem.
Every double zeta appearing in \eqref{tail:ordered-block} has even
outer order, odd inner order, and odd total weight. The classical
odd-weight reduction proved by Borwein, Borwein, and Girgensohn
\cite{BBG} gives the following explicit elimination. For even
\(a\ge2\), odd \(b\ge3\), and \(w=a+b\),
\begin{align}\label{tail:Euler-reduction}
 \zeta(a,b)={}&\frac12\left\{\binom wa-1\right\}\zeta(w)
              +\zeta(a)\zeta(b)\notag\\
 &-\sum_{k=1}^{(w-3)/2}
 \left\{\binom{w-2k-1}{b-1}+\binom{w-2k-1}{a-1}\right\}
                 \zeta(2k)\zeta(w-2k).
\end{align}
For the boundary case \(b=1\), Euler's formula is
\begin{equation}\label{tail:Euler-boundary}
 \zeta(a,1)=\frac a2\zeta(a+1)
            -\sum_{k=1}^{a/2-1}\zeta(2k)\zeta(a+1-2k).
\end{equation}
All zeta arguments on the right are at least two. These are classical
inputs, not novelty claims of this article.

\begin{corollary}[A moment-independent coordinate space]
\label{tail:fixed-space}
Let \(M=\max\{p,q\}\). Every value \(D_{p,q}(-2r)\), for
arbitrary \(r\ge0\), lies in the same rational span
\begin{equation}\label{tail:fixed-space-formula}
 \operatorname{span}_{\mathbb Q}
 \left(\{1\}\cup
       \{\zeta(2h):1\le h\le M/2-1\}\cup
       \{C_{p,q}(d):1\le d<M,\ d\text{ odd}\}\right).
\end{equation}
This is a spanning set of \(M\) specified numbers, each explicitly
reducible to ordinary zeta values and their pairwise products.
It is not a claim that those numbers are linearly independent.
\end{corollary}

\begin{proof}
For every odd \(d>M\), both ordered blocks in
\eqref{tail:ordered-block} contain only even zeta values of argument
at most \(M-2\), and the remaining block is rational. The odd
integers \(d<M\) give the finite list in
\eqref{tail:fixed-space-formula}. The moment formula proves the
inclusion, while \eqref{tail:Euler-reduction} and
\eqref{tail:Euler-boundary} give the stated ordinary-zeta reduction.
\end{proof}

This finiteness is uniform in the moment order. Increasing \(r\)
changes only rational Bernoulli coefficients; it introduces no new
arithmetic constants.

\subsection{The first mixed pair: a complete \texorpdfstring{$(2,4)$}{(2,4)} evaluation}

For \((p,q)=(2,4)\), write \(D_r=D_{2,4}(-2r)\). In this case
the entire block sequence has the short form
\begin{align}\label{tail:24-blocks}
 C_{2,4}(1)&=\frac{15}{2}\zeta(5)-3\zeta(2)\zeta(3),\notag\\
 C_{2,4}(3)&=\frac32\zeta(3)+\frac12\zeta(2),\notag\\
 C_{2,4}(2j+1)&=\frac{2j-1}{2}B_{2j-2}\zeta(2)
  -\frac{(2j-3)(2j^2-3j+7)}{24}B_{2j-4},\qquad j\ge2.
\end{align}
For example, the first line follows from
\(\zeta(2,3)=9\zeta(5)/2-2\zeta(2)\zeta(3)\) and
\(\zeta(4,1)=2\zeta(5)-\zeta(2)\zeta(3)\). Thus the arithmetic
space in this case can be written especially economically as
\begin{equation}\label{tail:24-space}
 \operatorname{span}_{\mathbb Q}
 \{1,\zeta(2),\zeta(3),5\zeta(5)-2\zeta(2)\zeta(3)\}.
\end{equation}

The first four values are
\begin{align}\label{tail:24-values}
 D_0={}&\frac{15}{2}\zeta(5)-3\zeta(2)\zeta(3),\notag\\
 D_1={}&-\frac58\zeta(5)+\frac14\zeta(2)\zeta(3)
                    +\frac12\zeta(3)+\frac16\zeta(2),\notag\\
 D_2={}&\frac7{32}\zeta(5)-\frac7{80}\zeta(2)\zeta(3)
              -\frac14\zeta(3)-\frac1{30}\zeta(2)-\frac3{40},\notag\\
 D_3={}&-\frac{155}{896}\zeta(5)+\frac{31}{448}\zeta(2)\zeta(3)
              +\frac7{32}\zeta(3)-\frac1{672}\zeta(2)+\frac{31}{672}.
\end{align}
The first two sums converge without subtraction. The next two yield
the polygamma identities
\begin{align}\label{tail:24-polygamma-sums}
 \sum_{n=0}^{\infty}
 \left\{\frac{x_n^4}{6}\psi'(n+1)\psi'''(n+1)-\frac13\right\}
       &=D_2,\notag\\
 \sum_{n=0}^{\infty}
 \left\{\frac{x_n^6}{6}\psi'(n+1)\psi'''(n+1)
                   -\frac{x_n^2}{3}+\frac7{36}\right\}
       &=D_3.
\end{align}
The signs and normalizations of both subtractions follow directly
from \(k_{2,4,0}=1/3\) and \(k_{2,4,1}=-7/36\).

There is also an exact finite-generation statement. The inverse
Bernoulli transform is
\begin{equation}\label{tail:inverse-transform}
 C_{p,q}(2j+1)=\sum_{r=0}^{j}
          \binom{2j+1}{2r}2^{-2(j-r)}D_{p,q}(-2r).
\end{equation}
It follows by multiplying the generating function in the next
section by \(2\sinh(t/2)/t\), or directly by the Bernoulli
generating identity. Put \(C_j=C_{2,4}(2j+1)\). Then
\[
 C_2=\frac14\zeta(2)-\frac38,
 \qquad C_3=-\frac1{12}\zeta(2)-\frac13,
\]
so
\[
 1=-\frac8{11}(C_2+3C_3),\qquad
 \zeta(2)=\frac{32C_2-36C_3}{11},\qquad
 \zeta(3)=\frac23C_1-\frac{32}{33}C_2+\frac{12}{11}C_3.
\]
Consequently \emph{every} \(D_r\) is an explicit rational linear
combination of \(D_0,D_1,D_2,D_3\). This is an exact generation
statement, independent of any conjectured arithmetic independence.
Two useful eliminations are
\begin{align}\label{tail:24-eliminations}
 7D_0+412D_1+2000D_2+1344D_3&=-88,\notag\\
 23D_0+524D_1-2480D_2-4032D_3&=176\zeta(2).
\end{align}
For instance the first is the single absolutely convergent identity
\begin{equation}\label{tail:rational-polygamma}
 \sum_{n=0}^{\infty}
 \left\{
 (1344x_n^6+2000x_n^4+412x_n^2+7)
       \frac{\psi'(n+1)\psi'''(n+1)}6
       -448x_n^2-\frac{1216}{3}
 \right\}=-88.
\end{equation}

